\Needspace{6\baselineskip}
\begin{center}\hypertarget{native-256-article183}{}\pdfbookmark[2]{183.}{article-183}183.\end{center}
Per methodum art. praec. mensura praecisionis pro ea sola incognita commode
exprimitur, cui in eliminationis negotio ultimus locus assignatus est,
quod incommodum ut evitemus, coëfficientem $\delta'''$ alio modo exprimere
conveniet. Ex aequationibus
\[
\begin{aligned}
P&=p'\\
Q&=q'+\frac{\beta}{\alpha}p'\\
R&=r'+\frac{\gamma'}{\beta'}q'+\frac{\gamma}{\alpha}p'\\
S&=s'+\frac{\delta''}{\gamma''}r'+\frac{\delta'}{\beta'}q'+\frac{\delta}{\alpha}p'
\end{aligned}
\]
sequitur, ipsas $p'$, $q'$, $r'$, $s'$ per $P$, $Q$, $R$, $S$ ita exprimi posse
\SourceLeaf{257}
\[
\begin{aligned}
p'&=P\\
q'&=Q+\mathfrak{A}P\\
r'&=R+\mathfrak{B}'Q+\mathfrak{A}'P\\
s'&=S+\mathfrak{C}''R+\mathfrak{B}''Q+\mathfrak{A}''P,
\end{aligned}
\]
ita ut $\mathfrak{A}$, $\mathfrak{A}'$, $\mathfrak{B}'$, $\mathfrak{A}''$,
$\mathfrak{B}''$, $\mathfrak{C}''$ sint quantitates determinatae. Erit itaque
(incognitarum numerum ad quatuor restringendo)
\[
s=-\frac{\lambda'''}{\delta'''}+\frac{\mathfrak{A}''}{\delta'''}P
+\frac{\mathfrak{B}''}{\delta'''}Q+\frac{\mathfrak{C}''}{\delta'''}R
+\frac{1}{\delta'''}S.
\]
Hinc conclusionem sequentem deducimus. Valores maxime probabiles incognitarum
$p$, $q$, $r$, $s$ etc. per eliminationem ex aequationibus $P=0$, $Q=0$,
$R=0$, $S=0$ etc. deducendi, manifesto, si aliquantisper $P$, $Q$, $R$, $S$
etc. tamquam indeterminatae spectentur, secundum eandem eliminationis
operationem in forma lineari per $P$, $Q$, $R$, $S$ etc. exprimentur,
ita ut habeatur
\[
\begin{gathered}
\begin{aligned}
p&=L+AP+BQ+CR+DS+\text{etc.}\\
q&=L'+A'P+B'Q+C'R+D'S+\text{etc.}\\
r&=L''+A''P+B''Q+C''R+D''S+\text{etc.}\\
s&=L'''+A'''P+B'''Q+C'''R+D'''S+\text{etc.}
\end{aligned}\\
\text{etc.}
\end{gathered}
\]
His ita factis, valores maxime probabiles ipsarum $p$, $q$, $r$, $s$ etc.
manifesto erunt resp. $L$, $L'$, $L''$, $L'''$ etc., mensuraque praecisionis
his determinationibus tribuendae resp. exprimetur per
$\BareRoot{\dfrac{1}{A}}\mathpunct{\raisebox{.45ex}{,}}$
$\BareRoot{\dfrac{1}{B'}}\mathpunct{\raisebox{.45ex}{,}}$
$\BareRoot{\dfrac{1}{C''}}\mathpunct{\raisebox{.45ex}{,}}$
$\BareRoot{\dfrac{1}{D'''}}$ etc., posita praecisione observationum
primitivarum $=1$. Quae enim de determinatione incognitae $s$ ante
demonstravimus (pro qua $\dfrac{1}{\delta'''}$ respondet ipsi $D'''$),
per solam incognitarum permutationem ad omnes reliquas transferre licebit.%
\HistoricalUnmarkedNote{[Handschriftliche Bemerkung:] Si gradus praecisionis
quantitatum $a$, $b$, $c$ etc. est $\alpha$, $\beta$, $\gamma$ etc.,
gradus praecisionis quantitatis $x=a+b+c+\text{etc.}$ erit
$\dfrac{1}{\BareRoot{\left(\dfrac{1}{\alpha\alpha}+\dfrac{1}{\beta\beta}+\dfrac{1}{\gamma\gamma}+\text{etc.}\right)}}\mathpunct{\cdot}$}
